\section{The information price of a functional}
\label{sec:infprice}

The gate needs a precision per persisted coordinate.  This section
prices it --- first for a single strike, then for the composite
objects a desk actually persists --- and meets, honestly, the one
place where the price is misquoted.

\subsection{Legs and baskets}

For the fitted volatility at one location $k_a$ --- call it a
\emph{leg} --- the natural currency is the quote support
$\mathcal{Q}_a=\mathcal{Q}(k_a)$ of \cref{eq:infsupport}: the
effective number of quotes sitting on that leg.  More effective
quotes, proportionally more precision; this is the usual
$1/n$ law of averaging $n$ measurements, taken as the working model.

But a desk does not think in isolated strikes.  Its vocabulary is
composite: the at-the-money level; the \emph{risk reversal}, the
volatility of a call-side strike minus that of a put-side strike ---
the trader's one-number summary of skew; the \emph{butterfly}, the
average of the two wing volatilities minus the level --- the
one-number summary of curvature.  Each is a \emph{signed basket} of
legs,
\begin{equation}
O=\sum_a \omega_a\,\sigma(k_a),
\label{eq:infbasket}
\end{equation}
with fixed coefficients $\omega_a$: the level has one leg with
$\omega=1$; the risk reversal has $\omega=(+1,-1)$ on its two wing
legs; the butterfly has
$\omega=(\tfrac12,\tfrac12,-1)$ on wing, wing, and money.  What is
the precision of a basket, given the support of its legs?

\begin{proposition}[The harmonic law]
\label{prop:infharmonic}
Suppose each leg volatility is estimable from today's quotes with
variance proportional to $1/\mathcal{Q}_a$, independently across
legs.  Then the implied estimate of the basket \cref{eq:infbasket}
has precision (in the same proportionality)
\begin{equation}
\mathcal{I}(O)
=\Bigl(\sum_a\frac{\omega_a^2}{\mathcal{Q}_a}\Bigr)^{\!-1}.
\label{eq:infharmonic}
\end{equation}
In particular, one unquoted leg ($\mathcal{Q}_a\to0$) drives
$\mathcal{I}(O)\to0$ regardless of how well the other legs are
quoted.
\end{proposition}

The reciprocal-of-summed-reciprocals shape is a (weighted) harmonic
mean of the legs' supports --- hence the name the chapter uses,
\emph{the harmonic law} --- and harmonic means are dominated by
their smallest entry, which is the dead-leg rule in one phrase.

\begin{proof}
For independent estimates, the variance of a linear combination is
the coefficient-squared-weighted sum of variances:
\[
\operatorname{Var}\Bigl(\sum_a\omega_a\hat\sigma_a\Bigr)
=\sum_a\omega_a^2\operatorname{Var}(\hat\sigma_a)
\propto\sum_a\frac{\omega_a^2}{\mathcal{Q}_a}.
\]
Precision is the reciprocal.  As any one $\mathcal{Q}_a\to0$, the
corresponding term $\omega_a^2/\mathcal{Q}_a\to\infty$, so the sum
diverges and $\mathcal{I}(O)\to0$
\citep[the standard sampling theory of a linear functional;][]{LehmannCasella1998}.
\end{proof}

The law has two consequences worth computing by hand.

\emph{The dead-leg rule.}  Take a risk reversal whose call leg is
well quoted, $\mathcal{Q}_{\mathrm{call}}=6$, but whose put leg has
almost nothing, $\mathcal{Q}_{\mathrm{put}}=0.1$.  Then
\[
\mathcal{I}(\mathrm{RR})
=\frac{1}{\tfrac{1}{6}+\tfrac{1}{0.1}}
=\frac{1}{0.167+10}
=\MacPriorBasketRrDead .
\]
Six quotes' worth of call-side evidence buys almost nothing: the
basket is a \emph{difference}, and a difference is only as known as
its least-known part.  At any requirement of order one, this basket's
gate stands essentially fully open --- correctly, because a risk
reversal with a missing put leg is a number the morning does not
know.

\emph{The factor of four.}  Now suppose both wing legs have the same
support $\mathcal{Q}$, while the money leg is quoted so densely that
its term in \cref{eq:infharmonic} contributes essentially nothing.
The risk reversal's coefficients are $\pm1$:
$\mathcal{I}(\mathrm{RR})=1/(2/\mathcal{Q})=\mathcal{Q}/2$.  The
butterfly's wing coefficients are $\tfrac12$, and they enter
\cref{eq:infharmonic} \emph{squared}:
$\mathcal{I}(\mathrm{BF})\approx1/(2\cdot\tfrac14/\mathcal{Q})
=2\mathcal{Q}$.  The same legs make the butterfly
\MacPriorBasketFactor\ times easier to identify than the risk
reversal.  Identifiability belongs to the functional, not to the
strikes: two baskets over identical legs can sit on opposite sides
of the requirement, and the gate --- which acts per functional ---
will treat them differently, exactly as it should.

\begin{figure}[htbp]
\centering
\paperfig{\textwidth}{fig_flt_basket.pdf}
\caption{The information price of a basket
(\cref{eq:infharmonic} on the support measure
\cref{eq:infsupport}).  (a)~A fixed budget of nine quotes spreads
over a widening band; the wing legs sit at
$\pm\MacPriorBasketLeg$ (dashed vertical).  The at-the-money leg is
identified from the first moment and slowly \emph{loses} precision as
the fixed budget thins; the butterfly crosses the dotted requirement
at half-width \MacPriorBasketCrossBf, the risk reversal only at
\MacPriorBasketCrossRr\ --- and both begin to rise before the band
actually reaches the legs, the kernel bleed discussed in the text.
(b)~The dead-leg rule: call-leg support fixed at six, put-leg support
swept to zero --- the basket's precision collapses with its weakest
leg, the butterfly running \MacPriorBasketFactor$\times$ above the
risk reversal throughout (its half-coefficients enter the harmonic
law squared, and the abundant money leg adds almost nothing).}
\label{fig:infbasket}
\end{figure}

\Cref{fig:infbasket} computes both effects.  In panel~(a), nine
quotes --- the size of the sparse morning cluster that returns in
\cref{sec:inffilter} --- are spread uniformly over a band whose
half-width grows along the horizontal axis, and the three baskets'
precisions are drawn against it.  The orange curve (the level) starts
high and \emph{declines}: a fixed budget spread ever wider leaves
fewer effective quotes at the money.  The green curve (butterfly) and
blue curve (risk reversal) start dark and rise as the band approaches
their $\pm\MacPriorBasketLeg$ legs, the butterfly crossing the
requirement first --- at half-width \MacPriorBasketCrossBf\ against
the risk reversal's \MacPriorBasketCrossRr\ --- and both roll over
and decline again once further widening only thins the legs.
Panel~(b) is the dead-leg rule as a curve: with the call leg held at
support six, the basket precisions fall to zero with the put leg's
support, the butterfly running \MacPriorBasketFactor$\times$ above
the risk reversal throughout.

\subsection{Where the price is misquoted}
\label{sec:infblind}

One honest caveat belongs here, because this is where it bites.
Support is a kernel sum, and a kernel radiates: quotes up to a
bandwidth away from a leg still credit it.  In panel~(a) both wing
baskets begin to rise --- and the butterfly actually crosses its
requirement --- \emph{before} the quoted band reaches the legs.  The
proxy declares a leg identified when the nearest quotes are still a
bandwidth short of it, which is precisely the region just beyond the
last quote that a wing prior exists to protect.  So the gate, fed by
support, can close on functionals the morning has not truly
measured.  The consequence is a design rule rather than a repair:
gated shape baskets cannot be the only mechanism holding a dark
wing, and \cref{sec:infcoords} pairs them with anchors placed by a
coverage test that does not radiate (the test is defined there).  The exact alternative ---
measure each functional's realized precision from a data-only fit,
then gate on that --- exists and costs one extra fit; the support
proxy is the price of gating before fitting.

We can now price any persisted functional and gate it.  What remains
open is the most consequential choice: \emph{which} functionals to
persist.  The next section shows that this choice --- mere
coordinates, seemingly --- decides whether the prior can damp a
market move.
